\subsection{Characteristic-prime curves and complete rank-two propagation}
\label{subsec:selmer-support-rank-two}

Keep the original rank-two normalized space, its arithmetic frame,
and every hypothesis of \texttt{Ok} and \texttt{Big}.
Write $A=A_{\mathscr S}$, $B=B_{\mathscr S}$,
$\Gamma=G_0\rtimes\langle c\rangle$, $d=\sum_vd_v$, and let
$\delta(c)=-1$.  The coefficient ring $\mathcal O$ is a complete
DVR with finite residue field $k$ of odd characteristic $\ell$ and
uniformizer $\varpi$; equal characteristic is allowed.  In particular
the original conditions retain
\begin{equation}\label{eq:ranktwo-original-budgets}
 \dim_k H^1_C(\Gamma,k(\Psi^{-2}\delta);0)\le C_0,
 \quad \dim_kH^1(\Gamma,k)\le C_0+1,
 \quad d>3C_0+7,\quad d_v>5.
\end{equation}
Here $H^i_C$ is the original continuous ordinary mapping-cone
cohomology, including its local trivializations.  We also retain the
original local inertia generators, the bounded exponent for unramified
anticyclotomic characters, the local separation $\Psi^2\ne1$, and
the Euler inequality with its exact no-cyclotomic-quotient hypothesis.
These are hypotheses about the actual frame, not consequences of
conductor one.\cite{support-ranktwo-source}

Let $\lambda_0\subset B$ be generated by $\varpi$ and
$\psi_{v,i}(g)-1$ for $g\in I_v$ with $\Psi(\operatorname{loc}_v g)=1$.
The original finite-weight condition on an ideal $I\subset B$ is
\begin{equation}\label{eq:ranktwo-finlam}
                  \mathfrak m_B\subseteq\sqrt{I+\lambda_0}.
\end{equation}
Whenever a prime containing $I$ occurs below, the inclusion is equality;
the inclusion formulation also covers the empty support locus.
Every chart below is the original local completion at a closed point
of the coefficient fibre with a nonzero lower entry, with its prescribed
finite unramified coefficient extension $\mathcal O'/\mathcal O$
and residue embedding.  Put $R=L.R$ and $k'=k(\mathcal O')$.

\begin{lemma}\label{lem:ranktwo-finlam-transfer}
For the actual ordinary Hecke ideal $J$, the original compatible map
$\Phi:B_{\rm det}\to B$ is surjective, and the original finite
weight-ring condition on $B_{\rm det}/J$ implies
\eqref{eq:ranktwo-finlam} for $I=\Phi(J)$.
\end{lemma}
\begin{proof}
The marked matrix difference is
$D=r(g_0)-r(g_0^*)=\operatorname{diag}(\zeta,-\zeta)$, with
$\zeta\bmod\mathfrak m=\Psi(g_0)-\Psi(g_0)^{-1}\ne0$.
The image $H$ of $B_{\rm det}$ in $\mathscr S.R$ is closed, since
the source is compact.  Undoing exactly the original $\kappa$ twist
and quotient map, it contains all traces and ordered local characters.
The identity
\[
 \zeta^2=\operatorname{tr}(r(g_0)^2-r(g_0)r(g_0^*)),\qquad
 e_1=\frac{D+\zeta}{2\zeta},\quad e_2=\frac{\zeta-D}{2\zeta}
\]
and Hensel's lemma put $\zeta$ in $H$.  Thus
$r(g)_{ii}=\operatorname{tr}(e_i r(g))$ and
$r(g)_{12}r(h)_{21}=\operatorname{tr}(e_1r(g)e_2r(h))$ lie in $H$.
Every degree-zero monomial is a product of diagonal entries and
opposite pairs; closure gives $B=H$.  Existence of $\Phi$ is the
determinant universal property applied to the same twisted pair.

If a prime $Q$ contains $I+\lambda_0$, its inverse image $Q'$ in
$B_{\rm det}$ contains $J$.  Modulo $Q$ the local characters equal
their residual values: raise inertia elements to $|k|-1$ and use the
simple roots of $X^{|k|-1}-1$.  The original weight universal property
therefore specializes the weight ring to its residue field.  Its
finite quotient $B_{\rm det}/Q'$ is a finite domain over that field,
hence a field.  Surjectivity of $\Phi$ makes $Q=\mathfrak m_B$.
The uniformizer in $\lambda_0$ is retained also in equal characteristic.
\end{proof}

\begin{theorem}[The actual rank-two curve lemma]
\label{thm:ranktwo-characteristic-curve}
Assume only the original frame conditions and
\eqref{eq:ranktwo-finlam}; no Hecke or modularity hypothesis is
needed for this assertion.  If $P\subset R$ contains $IR$ and
$\dim(R/P)\ge2d+1$, there is a prime $\mathfrak q\supset P$ with
\[
 \dim(R/\mathfrak q)=1,\qquad \varpi\in\mathfrak q,
 \qquad r_{\mathfrak q}\text{ absolutely irreducible},\qquad
 \psi_{v,1}|_{N_v}\ne\psi_{v,2}|_{N_v}\quad\text{for every }v,
\]
where $N_v=I_v\cap\ker\epsilon_v$ is the original subgroup.
This is an original \texttt{Nice} type-I prime.  In particular it
applies to the actual support of
Theorem~\ref{thm:selmer-support-starting-chart}.
\end{theorem}

\begin{proof}
\emph{The closed coefficient fibre.}
A finite power-series presentation of $\mathscr S.R$ has variables
of weights $0,2,-2$ and homogeneous relations: diagonal conjugation
preserves the original normalized equations.  Degree-zero monomials
are generated by the weight-zero variables and opposite pairs.
Dickson's lemma shows that each fixed homogeneous piece is a finite
$B$-module.  A fixed-weight monomial of sufficiently large total
degree contains arbitrarily many balanced pairs.  Consequently its
induced adic topology is its $\mathfrak m_B$-adic module topology.
These statements descend through the homogeneous relation ideal,
retain nilpotents, and show that entry functions modulo
$\mathfrak m_B$ are locally constant.

Choose the lower entry $u$ which is a unit in $R$.  Its product with
every upper entry is in $\mathfrak m_B$, so the fibre representation is
\[
 r(g)=\begin{pmatrix}\Psi(g)&0\\a(g)&\Psi(g)^{-1}\end{pmatrix},
 \quad a(gh)=a(g)\Psi(h)+\Psi(g)^{-1}a(h),\quad a(g^*)=a(g).
\]
The ordinary product and a local separating element give uniquely
$a(\operatorname{loc}_v h)=\gamma_v(\Psi(\operatorname{loc}_v h)
-\Psi(\operatorname{loc}_v h)^{-1})$.
The cocycle $z=a/\Psi$ extends to $\Gamma$ by $z(c)=0$, with
coefficient character $\Psi^{-2}\delta$ and these local
trivializations.  A coboundary with $z(c)=0$ is zero, since $2$ is
invertible.  Thus the strict cocycle space injects into the first
group of \eqref{eq:ranktwo-original-budgets}.  Finite residue
extension preserves its dimension over the new field.  Local
constancy and compactness let us expand all entries in a basis of
this space.  On $u\ne0$ the fibre algebra is a localization of an
algebra generated by at most $C_0$ coefficients.  Localization and
the selected local completion give
\begin{equation}\label{eq:ranktwo-fibre-dimension}
                         \dim R/\mathfrak m_BR\le C_0.
\end{equation}

\emph{The two bad loci.}
Let $T_-$ be generated by the upper entries.  Choose at most $C_0+1$
elements of $\Gamma$ detecting every homomorphism to
$\mathbb F_\ell$.  For their projections $t_j$ to $G_0$, use
$(t_jct_jc)^{|k|-1}$ and append the $d$ original local inertia
generators.  A finite $\ell$-group character killing these elements
has trivial $c$-invariant product, is unramified, and therefore has
exponent dividing the fixed integer in the original class-group
condition.  All chosen elements have $\Psi=1$.

Cut $R/T_-$ by $\varpi$ and the at most $d+C_0+1$ first diagonal
character values minus one.  In every prime quotient the residually
trivial diagonal character has that bounded exponent, by applying
the preceding assertion to every finite adic quotient.  A
characteristic-$\ell$ domain has no nontrivial $\ell$-power torsion,
and its prime-to-$\ell$ roots with residue one are one.  Both
diagonal characters are therefore residual; the original local
characteristic polynomials force the same for the ordered local
characters.  All remaining degree-zero monomials vanish or are
constant.  The cuts thus contain $\mathfrak m_B$ up to radical.
By \eqref{eq:ranktwo-fibre-dimension} and the height theorem,
\begin{equation}\label{eq:ranktwo-triangular-dimension}
                 \dim R/T_-\le d+2C_0+2<2d.
\end{equation}

For each $v$ let $C_v$ be generated by the differences of the two
ordered characters on $N_v$.  If a prime $Q\supset IR+\ell R+C_v$
is triangular, the preceding bound applies.  Otherwise both an upper
and a lower entry survive.  Every upper entry times the chosen lower
entry, and every lower entry times the chosen upper entry, belongs
to $B$.  For $\eta=Q\cap A$ and $\mathfrak p=\eta\cap B$ this gives
\[
 \operatorname{trdeg}_{\operatorname{Frac}(B/\mathfrak p)}
                 \operatorname{Frac}(A/\eta)\le1.
\]
This statement concerns the algebraized entry algebra.  To pass to
the completed chart, put $C=B/\mathfrak p$, $D=A/\eta$.
The point $x/\eta$ is closed over $\mathfrak m_C$ with finite
residue extension.  The complete Noetherian local ring $C$ is
universally catenary, so the dimension formula gives
\[
 \dim D_{x/\eta}=\dim C+
       \operatorname{trdeg}_{\operatorname{Frac}C}\operatorname{Frac}D.
\]
The selected local ring of $\mathcal O'\otimes_{\mathcal O}D$
is a localization of a finite etale extension of $D_{x/\eta}$;
the local map is flat with zero-dimensional closed fibre.
Its dimension is the same, completion preserves it, and $R/Q$ is
a quotient of that completion.  Thus, for every such $Q$,
\begin{equation}\label{eq:ranktwo-completed-dimension}
                         \dim R/Q\le\dim B/\mathfrak p+1.
\end{equation}
No transcendence-degree bound on a completed fraction field is used.
\cite{support-commutative-algebra}

At $v$, all but the one exceptional original inertia generator lie
in $N_v$; the two character coordinates there coincide modulo
$\mathfrak p$.  Retain both character coordinates at the other
places, and include $\varpi$.  There are at most $2d-d_v+2$ cuts.
The original generator property makes all residual-trivial inertial
characters one after these cuts.  Equation~\eqref{eq:ranktwo-finlam}
then leaves a zero-dimensional quotient.  Hence
$\dim B/\mathfrak p\le2d-d_v+2$, and
\begin{equation}\label{eq:ranktwo-coincidence-dimension}
                   \dim R/(IR+\ell R+C_v)<2d.
\end{equation}
The uniformizer cut is counted even when $\ell=0$ in $\mathcal O$.

\emph{Choosing the curve.}
A minimal prime over $P+(\varpi)$ has quotient dimension at least
$2d$, and so avoids $T_-$ and every $C_v$.  In that domain choose
a nonzero product $b$ of one element from each bad ideal.
While its dimension exceeds one, choose a maximal-ideal element
outside the finitely many height-one primes over $b$ and pass to a
height-one prime above it.  Such a prime avoids $b$.  Iteration gives
a one-dimensional prime avoiding all bad ideals; prime avoidance
works over the original finite residue field.  Pulling it back gives
$\mathfrak q$.  The two marked spectral projectors, a nonzero upper
entry and the unit lower entry generate both matrix units, proving
absolute irreducibility.  Avoiding $C_v$ is exactly the stated
ordered-character separation.  Normalization is finite and local
and gives a complete DVR $k''[[t]]$, without changing the kernel.
\end{proof}

\begin{theorem}[The original rank-two component theorem]
\label{thm:ranktwo-allgood}
Assume \eqref{eq:ranktwo-finlam} and the original condition
\[
 \mathfrak q\text{ Nice of type I},\ I\subset\mathfrak q
       \quad\Longrightarrow\quad
 I\subset P\quad\text{for every chart prime }P\subset\mathfrak q.
                                                               \tag{$P1'$}
\]
Assume also the original seed condition: for a stable prime
$\xi\subset A$ there is a prime $Z\subset B$ with
$I\subset Z\subset\xi\cap B$ and $\dim(B/Z)=1+2d$.
Then $I$ is contained in every stable prime of $A$.
This is the literal \texttt{AllGood} conclusion.
\end{theorem}

\begin{proof}
\emph{The actual chart functor.}
Its residual representation is a nonzero strict lower extension
$r_0$ as above.  Given a normalized lift over a complete local
$\mathcal O'$-algebra $C$ with the prescribed residue, let
$U\subset\mathfrak m_C$ be its upper-entry ideal and form
\begin{equation}\label{eq:ranktwo-laurent-rees}
 D_C=C[t^{-1},Ut]\subset C[t,t^{-1}],\qquad
 \mathfrak v=\mathfrak m_C+(\text{nonzero-degree elements}).
\end{equation}
Multiply the upper entries by $t$ and lower entries by $t^{-1}$.
Conjugation by $\operatorname{diag}(t,1)$ preserves the polarization
up to its common scalar, and preserves all original equations.
Balanced products give
$\mathfrak m_C^a\subset\mathfrak v^a\cap C
\subset\mathfrak m_C^{\lceil a/2\rceil}$.
Thus $C$ is closed in the vertex completion.  This remains true with
nilpotents: an element outside $\mathfrak v$ has nonzero reduction
as a Laurent polynomial over $k'$, so multiplication by it cannot
kill the first nonzero $\mathfrak m_C$-adic coefficient of another
Laurent polynomial.  Localization is injective; Krull intersection
then gives injectivity into the completion.

Base change the original weighted presentation to $\mathcal O'$.
Universality at the vertex sends $B$ into this closed $C$ and $A$
into $D_C$.  Evaluate only $D_C$ at $t=1$.  The resulting map
$A\to C$ has exactly the specified residue point, so the original
chart property gives $R\to C$.  Conversely such a map gives the
lift.  Entries and ordered characters prove uniqueness.  Hence this
is the deformation functor of the actual chart; no completed series
has been evaluated at $t=1$.

\emph{The presentation deficit.}
Let $W=M_2(k')$ with $G_0$ acting by $\operatorname{Ad}(r_0)$ and
$c:X\mapsto-J_2X^tJ_2$.  The local ordinary line is unique:
for a separating local element $s$, its projector is
\[
             E=\frac{r(s)-\psi_2(s)}{\psi_1(s)-\psi_2(s)}.
\]
The ordered product equations give
$r(g)E=\psi_1(g)E$ and
$(1-E)r(g)=\psi_2(g)(1-E)$.
Its preserving Lie algebra $\mathfrak b_v$ has dimension three.
The global invariants of $W$ are zero: the marking forces an
invariant to be diagonal, the nonzero lower extension forces it
to be scalar, and $c$ negates scalars.

The adjoint filtration has constituents
\[
                 \delta,\quad\Psi^{-2}\delta,\quad1,
                 \quad\Psi^2\delta.
\]
The original character exclusions \texttt{f2a}, \texttt{f2c} exclude
the cyclotomic character from this list.  Hence $W$ has no cyclotomic
quotient.  All hypotheses of the original Euler inequality apply,
including finite continuous residual image and the specified local
subspaces.  At each of the $d$ real places the fixed adjoint Lie
algebra is orthogonal of dimension one.  For the original cone,
\begin{equation}\label{eq:ranktwo-euler-deficit}
        h^2+3d\le h^1+d,\qquad h^1-h^2\ge2d.
\end{equation}

Strict equivalence classes of first-order ordinary deformations are
$H^1_C(\Gamma,W;(\mathfrak b_v))$.  Since $2$ is a unit,
a strict basis change normalizes the deformed symmetric form at $c$.
The remaining gauge Lie algebra is
$\{\operatorname{diag}(a,-a):a\in k'\}$, of dimension one,
with zero stabilizer.  The rank-two marking adds no equation:
$M-J_2M^tJ_2$ is already diagonal.  Thus
\begin{equation}\label{eq:ranktwo-tangent}
                       \dim\operatorname{Tan}(R/\mathcal O')=h^1+1.
\end{equation}
For a small extension, multiplication errors of matrix lifts and
errors of local line lifts are exactly a degree-two cocycle in the
same ordinary cone.  Changing lifts adds its coboundary.  Vanishing
of the class is equivalent to a correction; the symmetric-form
normalization restores $c$ afterwards.  Thus the full $H^2_C$ is
an obstruction space for this chart, including its local equations.

The dual of the minimal relation space injects into that obstruction
space.  Indeed a nonzero relation functional with zero obstruction
would split the corresponding universal small extension.  A change
of minimal tangent variables cannot change its selected relation
modulo the maximal ideal times the relation ideal.  This contradiction
also covers a relation involving the fixed coefficient uniformizer.
Therefore the actual chart has a presentation
\begin{equation}\label{eq:ranktwo-presentation}
 R=\mathcal O'[[X_1,\ldots,X_N]]/(f_1,\ldots,f_s),\qquad
 s\le h^2,\qquad N-s\ge2d+1.
\end{equation}
No regular-sequence assumption is used.

\emph{Local propagation.}
Grothendieck connectedness for the complete regular ring of dimension
$N+1$, with cohomological dimension at most $s$, makes the graph of
minimal primes connected when edges mean
\begin{equation}\label{eq:ranktwo-component-edges}
                         \dim R/(P+Q)\ge2d+1.
\end{equation}
Indeed a disconnected partition would disconnect the zero locus
after removing a subset of dimension less than $N-s$.
\cite{support-ranktwo-connectedness}
Every component has dimension at least $2d+2$ and is stable by
\eqref{eq:ranktwo-triangular-dimension}.  If $I\subset P$ and $Q$
is adjacent, a prime above $P+Q$ with dimension at least $2d+1$
gives a Nice curve by Theorem~\ref{thm:ranktwo-characteristic-curve}.
Condition $(P1')$ gives $I\subset Q$.  A starting prime containing
$I$ of dimension at least $2d+1$ gives the first good minimal prime
by the same curve argument.

The abstract seed suffices even with strict inclusion
$Z\subset\xi\cap B$.  Choose an upper entry $u$ and a lower entry
$v$ nonzero at $\xi$, and put $a=uv\in B\setminus Z$.  Exactly
\begin{equation}\label{eq:ranktwo-seed-localization}
                          A_a=B_a[u,u^{-1}].
\end{equation}
An upper entry $u'$ is $(u'v/a)u$ and a lower entry $v'$ is
$(uv')u^{-1}$; distinct degrees exclude a Laurent relation.
The prime $ZB_a[u,u^{-1}]$ contracts to a homogeneous stable prime
with coefficient contraction $Z$ and relative transcendence degree
one.  The graded Nakayama argument below places it below a closed
point of the coefficient-fibre open.  The dimension formula used
in \eqref{eq:ranktwo-completed-dimension}, with the selected etale
local extension and completion, gives a chart prime over it of
dimension $\dim(B/Z)+1=2d+2$.  This supplies a good initial chart.

\emph{Global propagation.}
Let $\mathrm{BB}$ be the original coefficient-fibre open with a
nonzero lower entry, and $S_0$ the finite-dimensional space of strict
cocycles in the fibre calculation.  There is a surjective continuous
map
\begin{equation}\label{eq:ranktwo-cocycle-fibre}
                         S_0\setminus\{0\}\longrightarrow\mathrm{BB}.
\end{equation}
A generic linear combination of a cocycle basis is a normalized
representation over its formal power-series ring.  The original
universal map restricts on $A$ to its polynomial coefficient family
and on $B$ to the residue map.  Conversely every field-valued point
of $\mathrm{BB}$ is precisely such a cocycle after scalar extension,
by local constancy and the fibre equations.  This proves the stated
surjectivity.  The seed makes the source nonempty, so $\mathrm{BB}$
is connected.

Every stable irreducible component meets $\mathrm{BB}$.  Its minimal
prime is homogeneous: its homogeneous core is prime by comparing
highest degrees of nonzero products, and minimality gives equality.
In the resulting graded domain, if the lower-entry ideal were in the
radical of the coefficient maximal ideal, sufficiently negative
homogeneous pieces would equal their maximal ideal multiples.
They are finite coefficient modules, so Nakayama would kill them,
contradicting the nonzero powers of a lower entry.  Thus there is a
lower-entry point in the closed fibre, and then a closed such point.

Call a component good if its minimal prime contains $I$.  At a closed
point of $\mathrm{BB}$, the selected coefficient extension and local
completion are faithfully flat.  If one component there is good,
local propagation makes every chart component good, and contraction
makes every original component there good.  All such components
are stable by \eqref{eq:ranktwo-presentation} and
\eqref{eq:ranktwo-triangular-dimension}.  The intersections of
$\mathrm{BB}$ with the good and remaining component unions are
therefore disjoint relatively closed sets: a nonempty intersection
would have a closed point, since the fibre is of finite type over
$k$.  They cover $\mathrm{BB}$, and the seed makes the good union
nonempty.  Connectedness gives that all stable components are good.
A stable prime contains a stable minimal prime, which proves the
full asserted prime quantifier.
\end{proof}

\paragraph{The arithmetic condition really gives $(P1')$.}
For $I=\Phi(J)$, normalize an original Nice curve to $k''[[t]]$,
twist the same representation and ordered characters by the original
$\kappa$, and descend through the same $\pi$.  The pair equations,
residual determinant and local flags are preserved.  The old inertia
equations follow from the original prime-to-$\ell$ exponent and the
same unit-resultant argument as in
Theorem~\ref{thm:selmer-support-seed-transfer}.  This is the actual
arithmetic point, with the same kernel.

On each original $N_v$, the ratio of its ordered normalized
characters is a nontrivial element of $1+tk''[[t]]$, a torsion-free
group.  Its powers enter any prescribed open subgroup while
retaining infinite order.  The exact local/cyclotomic identities of
the original twist put these elements in $G'=\ker\epsilon$.
This gives both original ordinary and cyclotomic-twisted inequalities
in \texttt{N1}.

For the full \texttt{N3} quantifier, take an open coset meeting $G'$
and its representative $\gamma\in G'$.  Choose a sufficiently small
open subgroup and such an $h$ in it with eigenvalue ratio
$\alpha/\beta$ of infinite order.  In an eigenbasis of $h$, write
$r(\gamma)=(m_{ij})$.  The discriminant of $r(\gamma)r(h)^a$,
up to a nonzero scalar, is the value at $z=(\alpha/\beta)^a$ of
\begin{equation}\label{eq:ranktwo-regular-coset}
 m_{11}^2z^2+(2m_{11}m_{22}-4\det r(\gamma))z+m_{22}^2.
\end{equation}
This polynomial is nonzero in odd characteristic: if the square
coefficients vanish, its middle coefficient is a unit.  Some power
therefore stays in the given coset with separable characteristic
polynomial.

The representation is type I on $U\cap G'$ for every open subgroup
$U$.  Replace $U$ by its normal core.  An irreducible representation
restricted to a normal subgroup is semisimple: the radical of its
finite-dimensional image algebra is stable under the full group and
must act by zero by irreducibility and nilpotence.  A reducible
rank-two restriction is either scalar, excluded by $h$, or a sum of
two distinct characters with an index-two line stabilizer $H$.
In the latter case there is a nontrivial quadratic self-twist $\xi$.
Each local ordinary line is an eigenline of a separating element
in $U\cap G'$, so every local ordinary group lies in $H$.
Reducing the self-twist and cancelling $\overline\kappa$ forces
$\xi=\Psi^2\circ\pi$.  This would make $\Psi^2$ trivial on
every local ordinary group, contrary to the original local
separation and surjectivity.  Thus the induced alternative is
impossible.  The argument persists under the finite coefficient
and fraction-field extensions used here.

Finally \texttt{N4} holds by the original modularity contract,
which quantifies over these arithmetic points.  The already proved
arithmetic P1 theorem, applied to the inverse image of any
$P\subset\mathfrak q$ under $B_{\rm det}\to R$, gives
$J\subset\Phi^{-1}(P)$, hence $I\subset P$.
Therefore Theorem~\ref{thm:ranktwo-allgood} applies to the actual
Hecke ideal.  Its conclusion is support propagation: the separate
rank-three type-E and thick-link argument, and the moving-family
integral Selmer estimate, are still required.
